\subsubsection{Evaluatiemetrieken voor phishingdetectie}
\label{subsubsec:evaluatiemetrieken}

\paragraph{Confusion matrix.}
Phishingdetectie wordt in de literatuur beschouwd als een binair classificatieprobleem, waarbij een e-mail als phishing of als legitiem wordt geclassificeerd. Door de inschatting te vergelijken met de werkelijkheid, kunnen vier mogelijke uitkomsten worden onderscheiden \autocite{Powers2020}:

\begin{itemize}
    \item \textbf{True positive (TP):} een phishingmail die als phishing wordt herkend;
    \item \textbf{True negative (TN):} een legitieme e-mail die als legitiem wordt herkend;
    \item \textbf{False positive (FP):} een legitieme e-mail die onterecht als phishing wordt ingeschat;
    \item \textbf{False negative (FN):} een phishingmail die als onterecht legitiem wordt ingeschat.
\end{itemize}

De uitkomsten worden gevisualiseerd in een confusion-matrix zoals in tabel \ref{tab:confusion-matrix}.

\begin{table}[htbp]
    \centering
    \footnotesize
    \setlength{\tabcolsep}{4pt}
    \begin{tabular}{lcc}
        \toprule
        & \begin{tabular}{c}\textbf{Ingeschat}\\\textbf{phishing}\end{tabular}
        & \begin{tabular}{c}\textbf{Ingeschat}\\\textbf{legitiem}\end{tabular} \\
        \midrule
        \textbf{Werkelijk phishing} & TP & FN \\
        \textbf{Werkelijk legitiem} & FP & TN \\
        \bottomrule
    \end{tabular}
    \caption{Confusion matrix voor phishingdetectie \autocite{Salloum2022}.}
    \label{tab:confusion-matrix}
\end{table}

Op basis van de aantallen uit deze matrix worden evaluatieformules afgeleid, zoals beschreven door \textcite{Sokolova2009}
\bigskip

\emph{Precision} meet de betrouwbaarheid van de positieve voorspellingen door te bepalen welk aandeel hiervan daadwerkelijk correct is:
\begin{equation*}
    \text{Precision} = \frac{TP}{TP + FP}
\end{equation*}

\emph{Recall} meet de volledigheid van de positieve voorspellingen door te bepalen welk aandeel van de daadwerkelijk positieve gevallen door het model is geïdentificeerd:
\begin{equation*}
    \text{Recall} = \frac{TP}{TP + FN}
\end{equation*}

Het \emph{false-positive-percentage} (FPR) toont welk deel van de legitieme e-mails ten onrechte als phishing wordt ingeschat:
\begin{equation*}
    \text{FPR} = \frac{FP}{FP + TN}
\end{equation*}

De \emph{F1-score} combineert precision en recall in één getal. Ze wordt berekend als het harmonisch gemiddelde van beide, waardoor ze enkel hoog is wanneer zowel precision als recall hoog zijn:
\begin{equation*}
    \text{F1-score} = 2 \cdot \frac{\text{Precision} \cdot \text{Recall}}{\text{Precision} + \text{Recall}}
\end{equation*}

Tot slot geeft \emph{accuracy} het aandeel correcte classificaties over alle e-mails:
\begin{equation*}
    \text{Accuracy} = \frac{TP + TN}{TP + TN + FP + FN}
\end{equation*}

In de praktijk worden deze maten niet even vaak gebruikt. In de analyse van \textcite{Salloum2022} was accuracy de meest gerapporteerde maat (83 experimenten), gevolgd door de F1-score (38), de confusion matrix (36) en precision en recall (elk 35) \autocite{Salloum2022}.
